\section{Proof of \Cref{thm:muPA-pi02-cons-muHA}}

% \subsection{Conservativity via double negation translation}
\label{sec:conservativity-via-double-negation}


\newcommand{\godelgentz}[1]{{#1}^{g}}
\newcommand{\friedman}[1]{{#1}^\rho}

In what follows we prove~\Cref{thm:muPA-pi02-cons-muHA}.
% \footnote{A different proof can be found in~\cite{Tupailo04doubleneg}.} 
% Henceforth, with little abuse of notation, with $\muPA$ (resp., $\muHA$) we denote the conservative extensions obtained by endowing the language  with new function symbols representing arbitrary primitive recursive functions, and by adding    equations defining the new identities as axioms. The following is a standard property:

% \begin{lemma}\label{lem:delta0-to-prim-rec}
%     Let  $\phi$ be quanfier-free. Then, there is a primitive recursive function symbol such that   $\vdash \phi(x,y) \leftrightarrow f(x,y)=0$, in both $\muPA$ and $\muHA$.
% \end{lemma}
% 
% 
We start with a definition a G\"{o}del-Gentzen-style translation of formulas of $\muPA$ into formulas of $\muHA$:
\[
\def\arraystretch{1.2}
\begin{array}{rcl}
    \godelgentz {(t=u)}  & \dfn & t=u \\
  \godelgentz {(t<u)}  & \dfn & \neg \neg t<u\\
  \godelgentz {(t \in X)} & \dfn & \neg \neg t \in X\\
  \godelgentz {(t \in \mu X\lambda x\phi)} & \dfn & \neg \neg t \in \mu X\lambda x\godelgentz{\phi} \\
\godelgentz {(\neg \phi)} & \dfn & \neg \godelgentz{\phi}\\
 \godelgentz {(\phi \wedge \psi)} & \dfn &  \godelgentz{\phi}\wedge \godelgentz{\psi}\\
  \godelgentz {(\phi \vee \psi)} & \dfn & \neg(\neg \godelgentz{\phi}\wedge \neg \godelgentz{\psi})\\
\godelgentz{(\forall x\phi)}  & \dfn & \forall x\godelgentz{\phi}\\
 \godelgentz{(\exists x\phi)} & \dfn & \neg (\forall x\neg \godelgentz{\phi})
\end{array}
\]

Some straightforward properties of the translation:

\begin{lemma}{\ }\label{lem:godel}
\begin{enumerate}
 \item \label{lem:godel1}$\godelgentz{(\phi(\psi))}= \godelgentz{\phi}(\godelgentz{\psi})$
    \item\label{lem:godel2} $\muHA \vdash \neg \neg \godelgentz {\phi} \to \godelgentz {\phi}$
\end{enumerate}
\end{lemma}

\begin{lemma}\label{lem:muPA-in-muHA}
    If $\muPA \vdash \phi$ then $\muHA \vdash \godelgentz{\phi}$.
\end{lemma}
\begin{proof}
We show that any axiom of $\muPA$ is derivable in $\muHA$, where the translation on equations is justified by the fact that $\muHA \vdash \neg \neg t=s \to t=s$. We just check the axioms for  the least fixed point operator $\mu$.  On the one hand, we have $ \godelgentz{\phi}(\mufo X x \godelgentz{\phi}, y) \limp y \in \mufo X x \godelgentz{\phi}$ as an instance of (pre),  and $y \in \mufo X x \godelgentz{\phi} \to \neg \neg y \in \mufo X x \godelgentz{\phi}$ by logic. We conclude by applying transitivity of implication and the generalisation rule introducing $\forall$. On the other hand, let us assume $\forall x(\godelgentz{\phi}(\godelgentz{\psi}, x) \limp \godelgentz{\psi}(x))$. By (ind) we obtain $y \in \mufo X x \godelgentz{\phi}  \limp \godelgentz{\psi}(y)$. By logic we have 
 $\neg \neg y \in \mufo X x \godelgentz{\phi}  \limp \neg \neg \godelgentz{\psi}(y)$ and $ \neg \neg \godelgentz{\psi}(y)\to \godelgentz{\psi}(y) $ (\Cref{lem:godel}.\ref{lem:godel2}). We conclude by applying transitivity of implication and deduction theorem, using~\Cref{lem:godel}.\ref{lem:godel1}. 
\end{proof}


We now use a standard trick called \emph{Friedman's translation} to infer our conservativity result from the above lemma. Let $\rho$ be a fixed formula of $\muHA$. For any formula $\phi$ of $\muHA$ we define $\friedman{\phi}$ as follows:
\[
\begin{array}{rcll}
    \friedman{(t=u)}  &  \dfn  & (t=u) \vee \rho\\
    \friedman{(t<u)}  &  \dfn  & (t<u) \vee \rho\\
    \friedman{(t \in X)}  &  \dfn  & (t \in X) \vee \rho\\
    \friedman{(t \in \mu X \lambda x\phi)}  &  \dfn  & (t \in \mu X\lambda x\friedman{\phi}) \vee \rho\\
   \friedman{(\neg \phi)}  & \dfn & \neg \friedman{\phi}\\
    \friedman{( \phi \circ \psi)}  & \dfn & ( \friedman{\phi}\circ \friedman{\psi})& \qquad \circ \in \{\wedge, \vee\}\\
    \friedman{\sigma x\phi} &\dfn & (\sigma x. \friedman{\phi}) &\qquad \sigma \in \{\forall, \exists\}
 \end{array}
\]
 The replacement must be done without variable capture, so whenever we consider, e.g., $\friedman{(\forall x \phi)}$, we must assume that $x$ is not free in $\rho$. 

Some straightforward properties of $\friedman{(\_)}$:
 \begin{lemma}{\ }\label{lem:friedman}
 \begin{enumerate}
     \item \label{lem:friedman1} $ \friedman{(\phi(\psi))}=\friedman{\phi}(\friedman{\psi})$
     \item \label{lem:friedman2}$\rho \vdash \friedman{\phi}$
 \end{enumerate}
 \end{lemma}

\begin{lemma}[Friedman's translation, adapted]
If $\muHA \vdash \phi$ then $\muHA \vdash \friedman{\phi}$.
\end{lemma}
\begin{proof}
       It is easy to check that $\Gamma \vdash \phi$ implies $\friedman{\Gamma}\vdash \friedman{\phi}$, where $\friedman{\Gamma}\dfn\{\friedman{\gamma} \ \vert \ \gamma \in \Gamma\}$, observing that $\rho \vdash \friedman{\phi}$ by Lemma~\ref{lem:friedman}.\ref{lem:friedman2}. So, Friedman's translation preserves derivability. Therefore, to conclude, it suffices to show that $\muHA$ proves Friedman's translation of $\muHA$'s axioms. We only consider the cases of the axioms for the least fixed point operator $\mu$.  On the one hand, by instantiating (pre), we have $\friedman{\phi}(\mufo X x \friedman{\phi}, y) \limp y \in \mufo X x \friedman{\phi}$, with $\friedman{\phi}(\mufo X x \friedman{\phi}, y)= \friedman{(\phi(\mufo X x\phi, y))}$ (Lemma~\ref{lem:friedman}.\ref{lem:friedman1}), so that we can conclude  $\friedman{\forall y (\phi(\mufo X x \phi, y) \limp y \in \mufo X x \phi)}$. On the other hand, we have to prove that $\forall x (\friedman{\phi}(\friedman{\psi}, x) \limp \friedman{\psi}(x)) $ implies $y \in \mufo X x \friedman{\phi} \vee \rho  \limp \friedman{\psi}(y)$.  By instantiating (ind) we have that the former implies $y \in \mufo X x \friedman{\phi}   \limp \friedman{\psi}(y)$. Also, from Lemma~\ref{lem:friedman}.\ref{lem:friedman2} we infer $\rho \to \friedman{\psi}(y)$, so $\forall x (\friedman{\phi}(\friedman{\psi}, x) \limp \friedman{\psi}(x))$ implies $(y \in \mufo X x \friedman{\phi}) \vee \rho  \limp \friedman{\psi}(y)$. This concludes the proof.
\end{proof}

We can finally prove Proposition~\ref{thm:muPA-pi02-cons-muHA}:

\begin{proof}[Proof of Proposition~\ref{thm:muPA-pi02-cons-muHA}]
Assume $\muPA \vdash \forall x\exists yA(x,y)$ for $A$ quantifier-free.  There is a primitive recursive function symbol $f$ such that $\vdash A(x,y)\leftrightarrow f(x,y)=0$ in both $\muPA$ and $\muHA$, so it suffices to show $\muHA \vdash \forall x\exists yf(x,y)=0$. We set $\rho\dfn \exists y .f(x,y)=0$.  Of course, $\muPA \vdash \rho$, so $\muHA \vdash \neg \neg \rho$ by Lemma~\ref{lem:muPA-in-muHA}. By Lemma~\ref{lem:friedman} we have  $\muHA \vdash (\friedman{\rho}\to \rho) \to \rho$, where $\friedman{\rho}= \exists y(f(x,y)=0 \vee \exists y.(x,y)=0)$, which is clearly equivalent to $\rho$. Thus, we have $\muHA \vdash (\rho\to \rho)\to \rho$, and so $\muHA\vdash \rho$. We conclude by generalising over $x$.
\end{proof}




       